\documentclass[conference]{IEEEtran}

\usepackage{amsmath,amssymb}
\usepackage{booktabs}
\usepackage{array}
\usepackage{url}

\newcolumntype{L}[1]{>{\raggedright\arraybackslash}p{#1}}

\newcommand{\method}{HKVM-RAG}
\newcommand{\whg}{Weighted-HG-KV}
\newcommand{\kgppr}{KG-PPR}

\begin{document}

\title{Supplementary Material for \method: Key-Value-Separated Hypergraph Evidence Organization for Multi-Hop RAG}

\author{
    \IEEEauthorblockN{Mingyu Zhang}
    \IEEEauthorblockA{\textit{Faculty of Computing} \\
    \textit{Harbin Institute of Technology} \\
    Harbin, China \\
    25b903142@stu.hit.edu.cn}
\and
    \IEEEauthorblockN{Ying Ma\textsuperscript{*}\thanks{\textsuperscript{*}Corresponding author.}}
    \IEEEauthorblockA{\textit{Faculty of Computing} \\
    \textit{Harbin Institute of Technology} \\
    Harbin, China \\
    y.ma@hit.edu.cn}
}

\maketitle

\section{Scope of the Supplementary Material}

This supplementary material contains reproducibility and diagnostic evidence that supports, but is not necessary to state, the main paper's claims. It follows the paper's bounded framing: \method{} is evaluated as a fixed-substrate evidence-indexing and dense-aware evidence-control method, not as a hidden-test benchmark, a standalone dense-retrieval replacement, or a full reproduction of every related RAG system. All tables below are derived from frozen prediction files and paper-facing exports. No new LLM extraction calls are introduced by these supplementary analyses.

\begin{table*}[t]
\caption{Supplementary evidence map. The supplement records statistical support, boundary diagnostics, and artifact provenance that would overburden the main paper; it does not introduce additional benchmark claims beyond the frozen result matrix.}
\label{tab:supp_evidence_map}
\centering
\scriptsize
\begin{tabular}{@{}L{0.18\textwidth}L{0.39\textwidth}L{0.36\textwidth}@{}}
\toprule
Supplement section & Supports & Does not support \\
\midrule
Source-level significance & WHG-KV-specific complementarity under the tested first-stage sources & Generic benefit from arbitrary structured sources \\
Controller significance & Dense-aware controller gains over ColBERTv2 and learned HKVM calibration & Hidden-test or deployment-throughput claims \\
Evidence-selection trajectory & Support-selection bottleneck and learned recovery path & Treating oracle rows as deployable systems \\
HotpotQA boundary & Separation of standalone failure and controller recovery & Standalone WHG-KV superiority on HotpotQA \\
Negative diagnostics & Need for learned control and cache-manifest discipline & Claim that LLM triples dominate all OpenIE alternatives \\
Artifact provenance & Audit path for paper-facing result families & Redistribution of non-redistributable provider or dataset artifacts \\
\bottomrule
\end{tabular}
\end{table*}

\section{Source-Level Significance}

Table~\ref{tab:supp_source_level_bootstrap} gives the paired-bootstrap support for the first-stage source diagnostics summarized in the main paper. The confidence interval tests WHG-KV against the corresponding first-stage baseline. The final column reports the point-estimate margin between WHG-KV and the strongest non-WHG structured complement in the same dataset--first-stage block.

\begin{table*}[t]
\caption{Source-level baseline-vs-WHG paired-bootstrap support and non-WHG margins. The confidence interval tests WHG-KV against the corresponding first-stage baseline; the final column is a point-estimate margin over the strongest matched non-WHG source. The table does not imply that all structured sources help or that the diagnostics cover all retriever families.}
\label{tab:supp_source_level_bootstrap}
\centering
\scriptsize
\begin{tabular*}{1.0\textwidth}{@{\extracolsep{\fill}} llrrr @{}}
\toprule
Dataset & First-stage & WHG-KV $\Delta$F1 vs base & 95\% CI & Margin over best non-WHG \\
\midrule
2WikiMultiHopQA & ColBERTv2 & $+11.084$ & $[+10.722,+11.442]$ & $+8.077$ \\
2WikiMultiHopQA & BM25 & $+16.452$ & $[+16.039,+16.865]$ & $+9.786$ \\
2WikiMultiHopQA & Contriever & $+23.119$ & $[+22.668,+23.577]$ & $+8.845$ \\
MuSiQue & ColBERTv2 & $+6.763$ & $[+5.813,+7.753]$ & $+8.219$ \\
MuSiQue & BM25 & $+12.038$ & $[+10.935,+13.120]$ & $+12.212$ \\
MuSiQue & Contriever & $+7.813$ & $[+6.810,+8.814]$ & $+7.605$ \\
HotpotQA & ColBERTv2 & $+5.966$ & $[+5.479,+6.456]$ & $+5.397$ \\
HotpotQA & BM25 & $+4.971$ & $[+4.554,+5.374]$ & $+4.188$ \\
HotpotQA & Contriever & $+5.166$ & $[+4.683,+5.634]$ & $+5.522$ \\
\bottomrule
\end{tabular*}
\end{table*}

The source-level exports use seeds 13, 29, and 43, 5,000 paper-export bootstrap samples per dataset--first-stage row, frozen train/dev prediction files, and no new LLM extraction calls. The artifact README records the ColBERTv2, BM25, Contriever, and HotpotQA source-level export roots, hashes, and redistributability notes.

\section{Dense-Aware Controller Significance}

The main paper reports paired-bootstrap F1 differences for the dense-aware HKVM controller against ColBERTv2 and learned HKVM calibration. Table~\ref{tab:supp_controller_bootstrap} complements the main paper by reporting the corresponding EM differences with 95\% confidence intervals. The controller improves EM over both ColBERTv2 and learned HKVM calibration on all three benchmarks. The 2Wiki gain over learned HKVM is intentionally small but positive, because learned HKVM calibration and the HKVM-only controller are already strong on that dataset.

\begin{table}[t]
\caption{Dense-aware controller EM significance. This table supplements the F1 bootstrap table in the main paper with exact-match (EM) paired-bootstrap results.}
\label{tab:supp_controller_bootstrap}
\centering
\scriptsize
\begin{tabular}{llrr}
\toprule
Dataset & Comparator & $\Delta$EM & 95\% CI \\
\midrule
2Wiki & ColBERTv2 & $+11.463$ & $[+11.116,+11.820]$ \\
2Wiki & Learned HKVM & $+0.477$ & $[+0.388,+0.565]$ \\
MuSiQue & ColBERTv2 & $+6.964$ & $[+5.984,+7.937]$ \\
MuSiQue & Learned HKVM & $+8.509$ & $[+7.509,+9.492]$ \\
HotpotQA & ColBERTv2 & $+5.991$ & $[+5.510,+6.482]$ \\
HotpotQA & Learned HKVM & $+2.921$ & $[+2.579,+3.255]$ \\
\bottomrule
\end{tabular}
\end{table}

The F1 counterparts of these intervals appear in the main paper's controller bootstrap table. Controller bootstrap exports use 5,000 samples over frozen prediction files with out-of-fold HKVM training signals. The source-level HotpotQA ColBERTv2 row in Table~\ref{tab:supp_source_level_bootstrap} is a separate source-replacement export, which explains its slightly different confidence interval from the main controller-vs-ColBERTv2 interval.

\section{Evidence-Selection Trajectory}

Table~\ref{tab:supp_trajectory} records the trajectory from no-gold WHG retrieval to oracle answer-path selection, train-derived support scoring, learned HKVM calibration, and the dense-aware controller. This table is a diagnostic trajectory rather than a single deployed pipeline. The oracle rows use support information and serve only to measure headroom.

\begin{table*}[t]
\caption{Evidence-selection trajectory from no-gold WHG retrieval to dense-aware control. The table supports the interpretation that support selection is a repairable bottleneck; oracle rows are diagnostic upper bounds and not deployable model results.}
\label{tab:supp_trajectory}
\centering
\scriptsize
\begin{tabular*}{1.0\textwidth}{@{\extracolsep{\fill}} l rrr rrr rrr @{}}
\toprule
& \multicolumn{3}{c}{2WikiMultiHopQA} & \multicolumn{3}{c}{MuSiQue} & \multicolumn{3}{c}{HotpotQA} \\
\cmidrule(lr){2-4} \cmidrule(lr){5-7} \cmidrule(lr){8-10}
Variant & F1 & EM & AR@10 & F1 & EM & AR@10 & F1 & EM & AR@10 \\
\midrule
No-gold WHG retrieval & 79.299 & 78.650 & 99.650 & 39.925 & 39.222 & 38.643 & 72.957 & 72.586 & 98.190 \\
Oracle answer-path selection & 90.884 & 90.832 & 99.634 & 74.330 & 74.059 & 73.066 & 91.111 & 91.047 & 98.190 \\
Train-derived support scorer & 87.212 & 86.967 & 98.378 & 54.252 & 53.662 & 50.228 & 81.525 & 81.269 & 96.016 \\
Learned HKVM calibration & 88.391 & 88.181 & 99.030 & 56.754 & 56.213 & 53.303 & 82.929 & 82.705 & 96.439 \\
Dense-aware HKVM controller & 88.846 & 88.658 & 100.000 & 65.073 & 64.722 & 64.115 & 85.810 & 85.627 & 100.000 \\
\bottomrule
\end{tabular*}
\end{table*}

The trajectory combines fixed-substrate dev outputs, oracle diagnostics, train-to-dev support-scoring outputs, learned HKVM calibration, and dense-aware controller exports. It should be read as a mechanism-oriented evidence path rather than as a single chronological production pipeline.

\section{Standalone Boundary on HotpotQA}

HotpotQA is included to prevent a standalone-retrieval overclaim. In the fixed-substrate standalone matrix, WHG-KV improves support coverage over KG-PPR but does not improve answer F1. Table~\ref{tab:supp_hotpot_boundary} separates that standalone boundary from the controller-mediated recovery reported in the main paper.

\begin{table}[t]
\caption{HotpotQA standalone boundary under matched candidate-passage budgets. The table blocks a standalone WHG-KV superiority claim while remaining compatible with controller-mediated recovery.}
\label{tab:supp_hotpot_boundary}
\centering
\scriptsize
\begin{tabular}{lrrr}
\toprule
Method & AR@10 & F1 & EM \\
\midrule
BM25 & 99.851 & 81.835 & 81.648 \\
Contriever & 100.000 & 79.930 & 79.730 \\
ColBERTv2 & 100.000 & 79.844 & 79.635 \\
KG-PPR & 95.381 & 73.645 & 73.315 \\
Static-HG & 98.190 & 73.340 & 73.018 \\
\whg{} & 98.190 & 72.957 & 72.586 \\
\bottomrule
\end{tabular}
\end{table}

This boundary is compatible with the dense-aware controller result: raw structured retrieval can underperform lexical and dense baselines on HotpotQA, while WHG-KV rank/score signals can still help a learned controller reorder answer-bearing passages.

\section{Negative and Boundary Diagnostics}

\subsection{Naive Fusion}

MuSiQue contains genuine ColBERTv2--WHG complementarity, but simple fixed-budget fusion and shallow single-feature gates did not reliably exceed ColBERTv2. RRF fusion of ColBERTv2 and HKVM predictions improved F1 over ColBERTv2 by only +0.110 on MuSiQue, far below the learned controller's +6.763. On 2WikiMultiHopQA and HotpotQA, RRF improved by +8.800 and +3.379 F1 respectively, but remained below the learned controller (+11.084 and +5.966). Three additional fixed-budget quota mixtures (dense3+hkvm2, dense4+hkvm1) were evaluated on MuSiQue and produced F1 values of 58.237, 58.309, and 57.923, none reliably exceeding ColBERTv2 (58.309). This motivates the learned passage-level controller used in the main paper, and blocks the simpler claim that generic RRF or quota fusion is enough.

\subsection{Controller Feature Ablation}

The dense-aware controller uses per-passage feature vectors containing ColBERTv2 and HKVM rank/score signals. Table~\ref{tab:supp_controller_ablation} decomposes the full feature set to identify which signal groups carry the controller gain. The ablation keeps the controller protocol, train-to-dev split, and hyperparameters fixed while varying the feature subset supplied to the logistic ranker.

\begin{table}[t]
\caption{Dense-aware controller feature ablation. Full uses both source signals plus explicit dense-HKVM interaction features; no-interaction drops the explicit interaction terms; score-only and rank-only keep only those feature groups; interaction-only uses only the explicit cross-source features. ColBERTv2 and learned HKVM calibration rows are reference predictors included in the main paper.}
\label{tab:supp_controller_ablation}
\centering
\scriptsize
\begin{tabular}{lrr}
\toprule
Variant & MuSiQue F1 & 2Wiki F1 \\
\midrule
ColBERTv2 (reference) & 58.309 & 77.763 \\
Learned HKVM calibration (reference) & 56.754 & 88.391 \\
\midrule
Controller dense-only & 58.247 & 77.679 \\
Controller HKVM-only & 56.760 & 88.415 \\
Controller interaction-only & 60.295 & 88.769 \\
Controller rank-only & 64.474 & 88.637 \\
Controller score-only & 65.254 & 88.696 \\
Controller no-interaction & 65.738 & 88.581 \\
Controller full & 65.073 & 88.846 \\
\bottomrule
\end{tabular}
\end{table}

The main finding is that explicit dense-HKVM interaction features are not necessary for the MuSiQue gain. The no-interaction variant reaches 65.738 F1 on MuSiQue, marginally exceeding the full controller (65.073). Score-only (65.254) also approaches full performance. Interaction-only (60.295) is far below full, confirming that interaction features alone are insufficient. On 2Wiki, HKVM calibration is already strong (88.391), so all dense+HKVM variants cluster within 0.7 F1 of each other. The safe mechanism claim is that dense and WHG-KV rank/score signals jointly carry useful passage-reordering information; explicit interaction features help on 2Wiki but are not the primary driver of the MuSiQue gain. This ablation uses out-of-fold HKVM training signals, seeds 13/29/43, 5,000-sample paired bootstrap, and no new LLM extraction calls.

\subsection{Extraction Substrate}

The shared LLM-triple cache is a controlled extraction substrate, not a claim that LLM triples dominate all OpenIE alternatives. In extraction-quality diagnostics, LLM triples were much sparser than heuristic OpenIE: objects per passage changed from 34.814 to 4.249 on 2Wiki and from 47.056 to 4.878 on MuSiQue. Pairwise KG-PPR with the sparse LLM tuple cache did not dominate heuristic OpenIE KG-PPR in that diagnostic. The paper therefore treats cache changes as new manifests rather than mixing them into the frozen result matrix.

\subsection{Hyperparameter Sensitivity}

Table~\ref{tab:supp_hyperparam_sensitivity} reports WHG-KV sensitivity to the three main hyperparameter groups: the confidence-weight coefficients $\lambda=(\lambda_f,\lambda_s,\lambda_b)$, the weight-scale parameter $\beta$, and the diffusion parameter $\rho$. Each variant varies one parameter while holding the others at their default values $(\lambda_f{=}0.20,\lambda_s{=}0.40,\lambda_b{=}0.40,\;\beta{=}2.0,\;\rho{=}0.35)$. All runs reuse the existing DeepSeek evidence tuple cache and evaluate on the MuSiQue dev answerable subset with seed 13.

\begin{table}[t]
\caption{WHG-KV hyperparameter sensitivity on MuSiQue (seed 13). Default row uses the paper-facing configuration. All variants reuse the frozen DeepSeek evidence tuple cache; F1 values are single-seed point estimates.}
\label{tab:supp_hyperparam_sensitivity}
\centering
\scriptsize
\begin{tabular}{llrrr}
\toprule
Group & Variant & AR@10 & F1 & EM \\
\midrule
\multicolumn{5}{@{}l}{\emph{Default}} \\
& $(\lambda_f{=}0.20,\lambda_s{=}0.40,\lambda_b{=}0.40),\;\beta{=}2.0,\;\rho{=}0.35$ & 37.815 & 38.946 & 38.229 \\
\midrule
\multicolumn{5}{@{}l}{\emph{$\beta$ (weight scale)}} \\
& $\beta=1.5$ & 37.360 & 38.693 & 37.940 \\
& $\beta=2.5$ & 38.395 & 39.482 & 38.767 \\
& $\beta=3.0$ & 38.395 & 39.834 & 39.139 \\
& $\beta=4.0$ & 39.015 & 39.918 & 39.222 \\
\midrule
\multicolumn{5}{@{}l}{\emph{$\rho$ (diffusion)}} \\
& $\rho=0.15$ & 38.477 & 39.383 & 38.726 \\
& $\rho=0.25$ & 37.857 & 39.292 & 38.602 \\
& $\rho=0.45$ & 37.650 & 38.937 & 38.146 \\
& $\rho=0.55$ & 36.947 & 38.471 & 37.650 \\
\midrule
\multicolumn{5}{@{}l}{\emph{$\lambda$ (confidence weights)}} \\
& Uniform $(0.33,0.33,0.33)$ & 37.940 & 38.988 & 38.271 \\
& Factual-only $(1,0,0)$ & 38.188 & 39.200 & 38.477 \\
& Salience-only $(0,1,0)$ & 37.898 & 38.936 & 38.229 \\
& Bridge-only $(0,0,1)$ & 37.940 & 39.022 & 38.271 \\
& Factual+Bridge $(0.5,0,0.5)$ & 38.022 & 38.943 & 38.229 \\
\bottomrule
\end{tabular}
\end{table}

Across all 14 variants, WHG-KV F1 spans 38.471--39.918 ($\Delta{=}1.447$), and all variants exceed the pairwise KG-PPR baseline (36.332 F1, Table~1 in the main paper). The default configuration is conservative: higher $\beta$ and lower $\rho$ modestly improve F1, but the paper's qualitative conclusions are robust to all tested hyperparameter perturbations. The default was chosen to avoid dev-set over-tuning rather than to maximize single-seed MuSiQue F1.

\section{Run and Artifact Provenance}

The public artifact is split across two hosting platforms:

\begin{itemize}
    \item \textbf{GitHub} (code, configs, scripts, results, manifests): \url{https://github.com/Mingyu-Zh/HKVM-RAG}
    \item \textbf{Hugging Face Dataset} (data, frozen outputs): \url{https://huggingface.co/datasets/MingY-Zh/HKVM-RAG}
\end{itemize}

This supplement itself is hosted in the GitHub repository at \texttt{supplemental\_material.pdf}. Reviewers should download \texttt{data/} and \texttt{frozen\_outputs/} from Hugging Face and place them at the repository root before running the paper-evidence verification script.

Table~\ref{tab:supp_artifacts} lists the paper-facing result families and the artifact records required to verify them. The artifact includes scripts (\texttt{scripts/verify\_paper\_evidence.py}, \texttt{scripts/reproduce\_fixed\_substrate.py}), configs (\texttt{configs/}), and a SHA-256 file manifest (\texttt{manifests/FILES.sha256}) for integrity verification.

\begin{table*}[t]
\caption{Paper-facing result provenance. The table identifies the protocol and artifact record for each result block; exact export roots, hashes, and redistributability constraints belong in the submitted artifact README.}
\label{tab:supp_artifacts}
\centering
\scriptsize
\setlength{\tabcolsep}{2.5pt}
\begin{tabular}{@{}L{0.16\textwidth}L{0.22\textwidth}L{0.19\textwidth}L{0.20\textwidth}L{0.17\textwidth}@{}}
\toprule
Result block & Verifies & Prediction protocol & Bootstrap/export & Artifact record \\
\midrule
P3 fixed-substrate matrix & Standalone WHG-vs-KG comparison and HotpotQA boundary & Frozen dev predictions; shared tuple cache; seeds 13/29/43 & Seed-13 1,000-sample WHG-minus-KG bootstrap; paper table export & \texttt{results/paper\_evidence/main\_fixed\_substrate/}, \texttt{frozen\_outputs/predictions/fixed\_substrate/} \\
Evidence-selection trajectory & Oracle headroom, support scorer, calibration, and controller sequence & Frozen dev predictions plus train-to-dev models; no dev labels at inference & Phase-specific paper exports; controller rows use 5,000-sample bootstrap & \texttt{results/paper\_evidence/calibration\_analysis/}, \texttt{results/run\_summaries/} \\
Phase1D controller & Dense-aware controller vs ColBERTv2 and learned HKVM & OOF train-side HKVM predictions; frozen dev predictions; seeds 13/29/43 & 5,000-sample paired bootstrap over frozen prediction files & \texttt{results/paper\_evidence/adaptive\_source\_controller/}, \texttt{frozen\_outputs/predictions/adaptive\_source\_controller/} \\
Source-level diagnostics & WHG-KV source specificity under ColBERTv2/BM25/Contriever & First-stage source fixed per block; structured source swapped; no new LLM calls & 5,000-sample baseline-vs-WHG bootstrap; non-WHG margins are point estimates & \texttt{results/paper\_evidence/source\_level\_ablation/}, bootstrap CSVs \\
Negative diagnostics & Naive-fusion and extraction-substrate boundaries & Diagnostic exports outside the frozen main matrix & Descriptive unless explicitly reported as paired bootstrap & \texttt{results/run\_summaries/extraction\_quality/}, prompt/schema records \\
\bottomrule
\end{tabular}
\end{table*}

\end{document}
